\documentclass[11pt]{article}
\usepackage[margin=1in]{geometry}
\usepackage{amsmath,amssymb,booktabs,hyperref,siunitx}
\title{Replication Report: Phase--Amplitude Collective-Mode Mixing in the\\ Imaginary CDW (Loop-Current) Phase of Kagome $A$V$_3$Sb$_5$}
\author{Automated replication (TEXTURES-100) of Xie \& Nagaosa, arXiv:2504.14166}
\date{Verdict: \textbf{REPLICATED}}
\begin{document}
\maketitle

\begin{abstract}
We reproduce from scratch the central theoretical claim of Xie \& Nagaosa
(arXiv:2504.14166): in the imaginary charge-density-wave (iCDW, loop-current)
phase of a triple-$Q$ kagome CDW, the amplitude (Higgs) and phase collective
modes \emph{mix} in the $A$ ($C_3$-symmetric) irrep, whereas in the real-CDW
(rCDW) reference they \emph{decouple}. Starting only from the mean-field free
energy (Eq.~2) and the fluctuation Lagrangian (Eq.~10), we minimize the free
energy to obtain the equilibrium order parameter and the symmetry-selected phase
$\theta_0$, build and diagonalize the $q=0$ $A$-channel dynamical matrix, and
recover the closed-form mixed-mode spectrum (Eq.~12) to machine precision
($\sim\!4\times10^{-16}$). The off-diagonal phase--amplitude susceptibility is
finite in the iCDW phase and vanishes (to $\sim\!10^{-17}$, i.e.\ exactly by
symmetry) in the rCDW phase, confirming the headline. A microscopic loop-current
cross-check using a shared kagome tight-binding kernel confirms the underlying
imaginary-hopping bond-current picture.
\end{abstract}

\section{Model and method}
The order parameter is the triple-$Q$ CDW matrix
$\Delta_{Q_j}(\mathbf r)=|\Delta_{Q_j}|e^{i(\mathbf Q_j\cdot\mathbf r+\theta_j)}$
on the kagome lattice (sublattices $A,B,C$; ordering vectors $Q_{1,2,3}$ at the
$M$ points). With $\lambda_3=0$ and imposed $C_3$ symmetry
($\theta_1=\theta_2=\theta_3\equiv\theta_0$, $|\Delta_{Q_1}|=|\Delta_{Q_2}|=|\Delta_{Q_3}|\equiv|\Delta_Q|$),
the mean-field free energy (paper Eq.~2) reduces to
\begin{equation}
F = 3(b\cos2\theta_0-\lambda_1')|\Delta_Q|^2 + \lambda_2|\Delta_Q|^3\cos3\theta_0
    + (3u_1+6u_2)|\Delta_Q|^4 .
\end{equation}
The dominant quadratic term selects $\theta_0=\pm\pi/2$ for $b>0$ (iCDW) and
$\theta_0=0,\pi$ for $b<0$ (rCDW). We minimize $F$ numerically over
$(|\Delta_Q|,\theta_0)$.

The $q=0$ $A$-channel fluctuation quadratic form, read directly from Eq.~(10) in
basis $(A^{(A)},\theta^{(A)})$, is
\begin{equation}
M=\begin{pmatrix}
4(u_1+2u_2)|\Delta_Q|^2 & \tfrac{3}{2}\lambda_2\sin(3\theta_0)|\Delta_Q|\\[2pt]
\tfrac{3}{2}\lambda_2\sin(3\theta_0)|\Delta_Q| & 2|b|-\tfrac{3}{2}\lambda_2\cos(3\theta_0)|\Delta_Q|
\end{pmatrix},
\end{equation}
and the mode energies satisfy $\det[M-\kappa_0\omega^2\mathbb{1}]=0$. The
off-diagonal element $\propto\sin(3\theta_0)$ is the mixing coefficient: nonzero
for iCDW ($\sin(3\theta_0)=\mp1$), exactly zero for rCDW ($\sin(3\theta_0)=0$).
For iCDW this yields the closed form (paper Eq.~12)
\begin{equation}
\kappa_0(\omega_\pm^{(A)})^2 = |b|+2(u_1+2u_2)|\Delta_Q|^2
   \pm\sqrt{\big(|b|-2(u_1+2u_2)|\Delta_Q|^2\big)^2+\tfrac{9}{4}\lambda_2^2|\Delta_Q|^2}.
\end{equation}

\section{Results}
Parameters (paper Fig.~3): $\lambda_1'=5,\ \lambda_2=0.1,\ u_1=u_2=0.5$;
$b=+1$ (iCDW), $b=-1$ (rCDW).

\begin{table}[h]\centering
\begin{tabular}{lcc}
\toprule
Quantity & iCDW ($b{=}+1$) & rCDW ($b{=}-1$)\\
\midrule
$\theta_0/\pi$ & $0.5$ & $1.0$\\
$|\Delta_Q|$ equilibrium & $1.41421$ & $1.42257$\\
$|\Delta_Q|$ analytic $\sqrt{(b+\lambda_1')/2(u_1{+}2u_2)}$ & $1.41421$ & --\\
off-diag mixing coeff & $-0.21213$ & $7.8\times10^{-17}$\\
cross susceptibility $\chi_{A\theta}$ & $8.86\times10^{-3}$ & $-2.9\times10^{-18}$\\
$\kappa_0\omega_-^2$ (numeric) & $1.99550$ & $2.21339$\\
$\kappa_0\omega_+^2$ (numeric) & $12.00450$ & $12.14226$\\
Eq.~(12) vs numeric max err & $4.4\times10^{-16}$ & --\\
modes mix? & \textbf{yes} & \textbf{no}\\
\bottomrule
\end{tabular}
\end{table}

\paragraph{Signature.} \verb|mixing_present_in_iCDW_only = true|;
\verb|Eq12_reproduced = true|. The numerically diagonalized $A$-channel spectrum
equals the paper's closed-form Eq.~(12) to machine precision, and the
phase--amplitude cross susceptibility is finite for iCDW and zero (by symmetry)
for rCDW.

\paragraph{Microscopic cross-check.} Using the shared kagome loop-current
mean-field kernel (\texttt{ollie\_loop\_current\_meanfield\_kernel}), the net
loop order at zero Peierls flux is $0$ (rCDW analog), while the loop-current
susceptibility to an imaginary bond flux is finite ($-62.2$), consistent with the
iCDW's imaginary inter-sublattice hopping generating bond currents.

\section{Verdict}
\textbf{REPLICATED.} The from-scratch mean-field + RPA-style fluctuation
calculation reproduces the paper's qualitative headline (mode mixing in iCDW,
decoupling in rCDW) \emph{and} the quantitative closed-form spectrum Eq.~(12) to
machine precision. Not attempted: the NV-center $T_1$ relaxometry / stray-field
magnitude estimate (Eqs.~13--15), which requires unspecified material/geometry
inputs.

\section*{Credits}
Microscopic loop-current cross-check uses the shared
\texttt{loop\_current\_meanfield\_kernel.py}
(\texttt{ollie\_loop\_current\_meanfield\_kernel}) from
\texttt{shared-kernels-cache/}. All numerical values are computed by
\texttt{report/evidence/xie2025\_replication.py}; results in
\texttt{work/xie2025\_result.json}.
\end{document}
